\section{Concurrency and Multithreading}

Modern computers have multiple processor cores, and C++ provides built-in support for concurrent programming starting with C++11. This section introduces the fundamentals of writing multithreaded C++ programs.

\subsection{Why Concurrency?}

Concurrency allows programs to:
\begin{itemize}
	\item \textbf{Improve performance}: Utilize multiple CPU cores for parallel computation
	\item \textbf{Maintain responsiveness}: Keep UIs responsive while performing background work
	\item \textbf{Handle I/O efficiently}: Process data while waiting for network or disk operations
	\item \textbf{Model real-world parallelism}: Some problems are naturally concurrent
\end{itemize}

\subsection{Creating Threads}

The \texttt{std::thread} class represents a thread of execution:

\begin{lstlisting}
	#include <thread>
	#include <iostream>
	
	void sayHello() {
		std::cout << "Hello from thread!\n";
	}
	
	int main() {
		std::thread t(sayHello);  // Create and start thread
		
		std::cout << "Hello from main!\n";
		
		t.join();  // Wait for thread to finish
		
		return 0;
	}
\end{lstlisting}

\subsubsection{Threads with Arguments}

\begin{lstlisting}
	void greet(const std::string& name, int times) {
		for (int i = 0; i < times; ++i) {
			std::cout << "Hello, " << name << "!\n";
		}
	}
	
	int main() {
		std::thread t(greet, "Alice", 3);
		t.join();
		return 0;
	}
\end{lstlisting}

\subsubsection{Threads with Lambdas}

\begin{lstlisting}
	int main() {
		int result = 0;
		
		std::thread t([&result]() {
			result = 42;
		});
		
		t.join();
		std::cout << "Result: " << result << "\n";
		
		return 0;
	}
\end{lstlisting}

\subsection{Joining and Detaching}

Every \texttt{std::thread} must be either joined or detached before destruction:

\begin{lstlisting}
	std::thread t(someFunction);
	
	// Option 1: Wait for thread to complete
	t.join();
	
	// Option 2: Let thread run independently
	t.detach();  // Thread continues even after t is destroyed
	
	// Destructor without join/detach calls std::terminate!
\end{lstlisting}

\begin{remark}
	Prefer \texttt{join()} over \texttt{detach()}. Detached threads are harder to manage and can cause issues if they access resources that have been destroyed.
\end{remark}

\subsection{Data Races and Synchronization}

When multiple threads access shared data, and at least one modifies it, you have a \textbf{data race}—undefined behavior:

\begin{lstlisting}
	int counter = 0;
	
	void increment() {
		for (int i = 0; i < 100000; ++i) {
			++counter;  // DATA RACE! Multiple threads modifying counter
		}
	}
	
	int main() {
		std::thread t1(increment);
		std::thread t2(increment);
		
		t1.join();
		t2.join();
		
		// counter is NOT guaranteed to be 200000!
		std::cout << counter << "\n";  // Could be any value
	}
\end{lstlisting}

\subsection{Mutexes}

A \texttt{std::mutex} provides mutual exclusion—only one thread can hold it at a time:

\begin{lstlisting}
	#include <mutex>
	
	int counter = 0;
	std::mutex mtx;
	
	void increment() {
		for (int i = 0; i < 100000; ++i) {
			mtx.lock();
			++counter;      // Protected by mutex
			mtx.unlock();
		}
	}
\end{lstlisting}

\subsubsection{std::lock\_guard (RAII Locking)}

Never use raw \texttt{lock()}/\texttt{unlock()}—use RAII wrappers:

\begin{lstlisting}
	void increment() {
		for (int i = 0; i < 100000; ++i) {
			std::lock_guard<std::mutex> lock(mtx);  // Locks on construction
			++counter;
			// Automatically unlocks when lock goes out of scope
		}
	}
\end{lstlisting}

\subsubsection{std::unique\_lock (Flexible Locking)}

\texttt{std::unique\_lock} offers more flexibility than \texttt{lock\_guard}:

\begin{lstlisting}
	std::unique_lock<std::mutex> lock(mtx);
	
	// Can unlock and relock
	lock.unlock();
	// ... do work that doesn't need the lock ...
	lock.lock();
	
	// Can transfer ownership
	std::unique_lock<std::mutex> lock2 = std::move(lock);
	
	// Can defer locking
	std::unique_lock<std::mutex> lock3(mtx, std::defer_lock);
	// ... later ...
	lock3.lock();
\end{lstlisting}

\subsection{Deadlocks}

A \textbf{deadlock} occurs when threads wait for each other indefinitely:

\begin{lstlisting}
	std::mutex mtx1, mtx2;
	
	void thread1() {
		std::lock_guard<std::mutex> lock1(mtx1);
		std::this_thread::sleep_for(std::chrono::milliseconds(1));
		std::lock_guard<std::mutex> lock2(mtx2);  // Waits for mtx2
	}
	
	void thread2() {
		std::lock_guard<std::mutex> lock2(mtx2);
		std::this_thread::sleep_for(std::chrono::milliseconds(1));
		std::lock_guard<std::mutex> lock1(mtx1);  // Waits for mtx1
	}
	// DEADLOCK: thread1 holds mtx1, waits for mtx2
	//           thread2 holds mtx2, waits for mtx1
\end{lstlisting}

\subsubsection{Avoiding Deadlocks}

Use \texttt{std::lock} to lock multiple mutexes atomically:

\begin{lstlisting}
	void safeFunction() {
		std::unique_lock<std::mutex> lock1(mtx1, std::defer_lock);
		std::unique_lock<std::mutex> lock2(mtx2, std::defer_lock);
		
		std::lock(lock1, lock2);  // Locks both without deadlock
		
		// ... work with protected data ...
	}
	
	// C++17: std::scoped_lock handles multiple mutexes automatically
	void saferFunction() {
		std::scoped_lock lock(mtx1, mtx2);  // Locks both, unlocks on destruction
		// ... work with protected data ...
	}
\end{lstlisting}

\subsection{Condition Variables}

\texttt{std::condition\_variable} allows threads to wait for conditions:

\begin{lstlisting}
	#include <condition_variable>
	#include <queue>
	
	std::queue<int> dataQueue;
	std::mutex mtx;
	std::condition_variable cv;
	
	void producer() {
		for (int i = 0; i < 10; ++i) {
			{
				std::lock_guard<std::mutex> lock(mtx);
				dataQueue.push(i);
			}
			cv.notify_one();  // Wake up one waiting consumer
		}
	}
	
	void consumer() {
		while (true) {
			std::unique_lock<std::mutex> lock(mtx);
			
			cv.wait(lock, []{ return !dataQueue.empty(); });  // Wait for data
			
			int value = dataQueue.front();
			dataQueue.pop();
			
			lock.unlock();
			
			std::cout << "Consumed: " << value << "\n";
		}
	}
\end{lstlisting}

\subsection{Atomic Operations}

For simple operations, \texttt{std::atomic} provides lock-free thread safety:

\begin{lstlisting}
	#include <atomic>
	
	std::atomic<int> counter{0};
	
	void increment() {
		for (int i = 0; i < 100000; ++i) {
			++counter;  // Atomic increment - no mutex needed!
		}
	}
	
	int main() {
		std::thread t1(increment);
		std::thread t2(increment);
		
		t1.join();
		t2.join();
		
		std::cout << counter << "\n";  // Guaranteed to be 200000
	}
\end{lstlisting}

\subsection{Futures and Promises}

\texttt{std::future} and \texttt{std::promise} provide a way to transfer values between threads:

\begin{lstlisting}
	#include <future>
	
	int computeValue() {
		std::this_thread::sleep_for(std::chrono::seconds(2));
		return 42;
	}
	
	int main() {
		// std::async runs function asynchronously, returns future
		std::future<int> result = std::async(std::launch::async, computeValue);
		
		std::cout << "Doing other work...\n";
		
		// get() blocks until result is ready
		std::cout << "Result: " << result.get() << "\n";
		
		return 0;
	}
\end{lstlisting}

\subsubsection{Promise and Future Pair}

\begin{lstlisting}
	void producer(std::promise<int>& prom) {
		// Do some work...
		prom.set_value(42);  // Fulfill the promise
	}
	
	int main() {
		std::promise<int> prom;
		std::future<int> fut = prom.get_future();
		
		std::thread t(producer, std::ref(prom));
		
		std::cout << "Waiting for result...\n";
		std::cout << "Got: " << fut.get() << "\n";
		
		t.join();
		return 0;
	}
\end{lstlisting}

\subsection{Thread-Local Storage}

\texttt{thread\_local} variables have separate instances for each thread:

\begin{lstlisting}
	thread_local int threadId = 0;
	
	void worker(int id) {
		threadId = id;  // Each thread has its own copy
		
		// ... later ...
		std::cout << "Thread " << threadId << " working\n";
	}
\end{lstlisting}

\subsection{Best Practices}

\begin{enumerate}
	\item \textbf{Minimize shared mutable state}: The less data threads share, the fewer synchronization issues
	
	\item \textbf{Use high-level constructs}: Prefer \texttt{std::async}, \texttt{std::future} over raw threads when possible
	
	\item \textbf{Always use RAII for locks}: \texttt{lock\_guard}, \texttt{unique\_lock}, \texttt{scoped\_lock}
	
	\item \textbf{Prefer atomics for simple operations}: Faster than mutexes for counters and flags
	
	\item \textbf{Be aware of false sharing}: Keep frequently-modified atomic variables on separate cache lines
	
	\item \textbf{Test thoroughly}: Concurrency bugs are notoriously hard to reproduce and debug
\end{enumerate}

\begin{remark}
	Concurrent programming is challenging. Start with simple patterns, understand the fundamentals thoroughly, and consider using higher-level libraries (like Intel TBB or cpp-taskflow) for complex scenarios.
\end{remark}